% ========================================================
% BÀI TẬP TOÁN 5 - KẾT NỐI TRI THỨC VỚI CUỘC SỐNG
% Bài 70: Ôn tập tỉ số, tỉ số phần trăm
% Định dạng: BTTeX Part III - Trắc nghiệm điền số ngắn
% ========================================================

% Câu 1
\begin{ex}
  Tìm hai số có tổng là 100 và tỉ số là $\frac{2}{3}$. Số bé là bao nhiêu?
  \shortans{40}
  \loigiai{
    Tổng số phần: $2 + 3 = 5$. Giá trị 1 phần: $100 : 5 = 20$. Số bé: $20 \times 2 = 40$.
  }
\end{ex}

% Câu 2
\begin{ex}
  Ở bài trên, số lớn là bao nhiêu?
  \shortans{60}
  \loigiai{
    $20 \times 3 = 60$.
  }
\end{ex}

% Câu 3
\begin{ex}
  Tìm hai số có hiệu là 45 và tỉ số là $\frac{4}{7}$. Số bé là bao nhiêu?
  \shortans{60}
  \loigiai{
    Hiệu số phần: $7 - 4 = 3$. Giá trị 1 phần: $45 : 3 = 15$. Số bé: $15 \times 4 = 60$.
  }
\end{ex}

% Câu 4
\begin{ex}
  Ở bài trên, số lớn là bao nhiêu?
  \shortans{105}
  \loigiai{
    $15 \times 7 = 105$.
  }
\end{ex}

% Câu 5
\begin{ex}
  Một lớp có 30 học sinh, trong đó có 12 bạn nam. Tỉ số phần trăm của số bạn nam là bao nhiêu phần trăm?
  \shortans{40}
  \loigiai{
    $(12 : 30) \times 100\% = 40\%$.
  }
\end{ex}

% Câu 6
\begin{ex}
  Ở lớp học trên, tỉ số phần trăm số bạn nữ là bao nhiêu phần trăm?
  \shortans{60}
  \loigiai{
    $100\% - 40\% = 60\%$.
  }
\end{ex}

% Câu 7
\begin{ex}
  Tìm $15\%$ của 200. Kết quả là bao nhiêu?
  \shortans{30}
  \loigiai{
    $200 \times 15\% = 30$.
  }
\end{ex}

% Câu 8
\begin{ex}
  Tìm một số biết $25\%$ của nó là 50. Số đó là bao nhiêu?
  \shortans{200}
  \loigiai{
    $50 : 25\% = 200$.
  }
\end{ex}

% Câu 9
\begin{ex}
  Một cửa hàng bán chiếc xe đạp giá 2 triệu đồng, nhân dịp lễ giảm giá $10\%$. Số tiền được giảm là bao nhiêu nghìn đồng?
  \shortans{200}
  \loigiai{
    $2\,000 \times 10\% = 200$ nghìn đồng.
  }
\end{ex}

% Câu 10
\begin{ex}
  Ở bài trên, giá bán sau khi giảm là bao nhiêu trăm nghìn đồng?
  \shortans{18}
  \loigiai{
    $2\,000 - 200 = 1\,800$ nghìn đồng = 18 trăm nghìn đồng.
  }
\end{ex}

% Câu 11
\begin{ex}
  Một người gửi tiết kiệm 50 triệu đồng với lãi suất $6\%$ một năm. Sau 1 năm tiền lãi nhận được là bao nhiêu triệu đồng?
  \shortans{3}
  \loigiai{
    $50 \times 6\% = 3$ triệu đồng.
  }
\end{ex}

% Câu 12
\begin{ex}
  Dân số một xã là $10\,000$ người, tỉ lệ tăng dân số hàng năm là $1{,}2\%$. Sau một năm dân số xã tăng thêm bao nhiêu người?
  \shortans{120}
  \loigiai{
    $10\,000 \times 1{,}2\% = 120$ người.
  }
\end{ex}

% Câu 13
\begin{ex}
  Một mảnh đất hình chữ nhật có chu vi $160\text{ m}$, chiều rộng bằng $\frac{3}{5}$ chiều dài. Nửa chu vi là bao nhiêu mét?
  \shortans{80}
  \loigiai{
    $160 : 2 = 80\text{ m}$.
  }
\end{ex}

% Câu 14
\begin{ex}
  Ở bài trên, chiều rộng của mảnh đất là bao nhiêu mét?
  \shortans{30}
  \loigiai{
    Tổng số phần: $3 + 5 = 8$. Chiều rộng: $(80 : 8) \times 3 = 30\text{ m}$.
  }
\end{ex}

% Câu 15
\begin{ex}
  Ở bài trên, chiều dài của mảnh đất là bao nhiêu mét?
  \shortans{50}
  \loigiai{
    $(80 : 8) \times 5 = 50\text{ m}$.
  }
\end{ex}

% Câu 16
\begin{ex}
  Ở bài trên, diện tích của mảnh đất là bao nhiêu mét vuông?
  \shortans{1500}
  \loigiai{
    $50 \times 30 = 1\,500\text{ m}^2$.
  }
\end{ex}

% Câu 17
\begin{ex}
  Hiện nay mẹ hơn con 28 tuổi. Tuổi con bằng $\frac{2}{9}$ tuổi mẹ. Tuổi con hiện nay là bao nhiêu tuổi?
  \shortans{8}
  \loigiai{
    Hiệu số phần: $9 - 2 = 7$. Tuổi con: $(28 : 7) \times 2 = 8$ tuổi.
  }
\end{ex}

% Câu 18
\begin{ex}
  Ở bài trên, tuổi mẹ hiện nay là bao nhiêu tuổi?
  \shortans{36}
  \loigiai{
    $8 + 28 = 36$ tuổi.
  }
\end{ex}

% Câu 19
\begin{ex}
  Tỉ số học sinh nam và nữ của khối 5 là $\frac{4}{5}$. Biết số học sinh nữ nhiều hơn nam là 20 bạn. Khối 5 có tất cả bao nhiêu học sinh?
  \shortans{180}
  \loigiai{
    Hiệu số phần: $5 - 4 = 1$ phần = 20 bạn. Tổng số phần: $4 + 5 = 9$. Tổng học sinh: $20 \times 9 = 180$ bạn.
  }
\end{ex}

% Câu 20
\begin{ex}
  Một tổ dự định làm 400 sản phẩm, đã làm được 480 sản phẩm. Tổ đó đã hoàn thành bao nhiêu phần trăm kế hoạch?
  \shortans{120}
  \loigiai{
    $(480 : 400) \times 100\% = 120\%$.
  }
\end{ex}

% Câu 21
\begin{ex}
  Ở bài trên, tổ đó đã vượt mức kế hoạch bao nhiêu phần trăm?
  \shortans{20}
  \loigiai{
    $120\% - 100\% = 20\%$.
  }
\end{ex}

% Câu 22
\begin{ex}
  Một bể nuôi cá có 30 con cá vàng và 20 con cá bảy màu. Tỉ số phần trăm cá vàng trong bể là bao nhiêu phần trăm?
  \shortans{60}
  \loigiai{
    Tổng số cá: 50 con. $(30 : 50) \times 100\% = 60\%$.
  }
\end{ex}

% Câu 23
\begin{ex}
  Một chiếc tivi có giá 8 triệu đồng, được giảm giá $15\%$. Khách hàng phải trả bao nhiêu trăm nghìn đồng?
  \shortans{68}
  \loigiai{
    Giảm: $8 \times 15\% = 1{,}2$ triệu. Còn lại: $8 - 1{,}2 = 6{,}8$ triệu = 68 trăm nghìn.
  }
\end{ex}

% Câu 24
\begin{ex}
  Tìm $x$ biết: $x : 15 = 4 : 5$. Giá trị của $x$ là bao nhiêu?
  \shortans{12}
  \loigiai{
    $x = (15 \times 4) : 5 = 12$.
  }
\end{ex}

% Câu 25
\begin{ex}
  Số học sinh giỏi của trường là 240 em, chiếm $30\%$ số học sinh toàn trường. Trường đó có bao nhiêu học sinh?
  \shortans{800}
  \loigiai{
    $240 : 30\% = 800$ học sinh.
  }
\end{ex}
